\section{ICA 与 CCA：选择不同的可辨认结构}
\label{sec:13-Machine-Learning/04-Dimensionality-Reduction/07}

一段 64 通道的脑电记录，每次眨眼都在所有通道上留下一道大幅波形。你想把眨眼伪迹去掉，保留皮层信号。第一反应是 PCA：伪迹幅度大，应该落在前几个主成分里，减掉就行。试了一遍，波形是干净了些，可几个明显的诱发响应也一起变弱了——被减掉的成分里既有眨眼，也有真实活动。上一节刚讲过原因：似然（或方差）只能定住子空间，定不住子空间里的具体方向，而你要减掉的偏偏是一个具体方向。这一节讲两种把方向重新钉住的办法，代价是各自多押一条假设。

\subsection{白化之后还剩一个旋转，二阶统计量交不出答案}

把问题写成标准形式。源信号 \(s\in\R^k\)，观测是它们的线性混合，

\[
x = As .
\]

盲源分离要只凭许多时刻的 \(x\) 把 \(s\) 和 \(A\) 都找回来。先做能做的：中心化，再白化，即找线性变换使 \(\Cov(\tilde x)=I\)。这一步用掉了协方差里的全部信息，也确实有用——它把混合矩阵约化到只差一个正交阵。

问题正在这里。若 \(\Cov(\tilde x)=I\)，那么对任意正交 \(R\)，\(\Cov(R\tilde x)=RR^{\trans}=I\) 同样成立。二阶统计量对旋转完全不敏感，PCA 也好、因子分析也好，都只用到二阶统计量，它们没有任何依据从这一族旋转里挑出真正的源方向。要往前走，必须动用更高阶的信息。

\begin{center}
\begin{tikzpicture}[>=stealth,scale=1.0]
% 面板一：非高斯源
\begin{scope}[rotate=32]
\draw[dashed,black!40] (-1.25,-1.25) rectangle (1.25,1.25);
\foreach \a in {-1.1,-0.55,0,0.55,1.1}
  \foreach \b in {-1.1,-0.55,0,0.55,1.1}
    \fill[black!72] (\a,\b) circle (1.5pt);
\foreach \a in {-0.8,-0.25,0.3,0.85}
  \foreach \b in {-0.8,-0.25,0.3,0.85}
    \fill[black!72] (\a,\b) circle (1.5pt);
\draw[very thick,->] (0,0) -- (1.75,0);
\draw[very thick,->] (0,0) -- (0,1.75);
\end{scope}
\draw[dashed,black!55,->] (0,0) -- (1.9,0);
\draw[dashed,black!55,->] (0,0) -- (0,1.9);
\node[font=\scriptsize,anchor=west] at (1.55,1.75) {源方向};
\node[font=\scriptsize,anchor=west] at (1.95,0.12) {白化后任取的正交轴};
\node[font=\small] at (0,-2.4) {源非高斯：边缘暴露了旋转};

% 面板二：高斯源
\begin{scope}[shift={(6.6,0)}]
\draw[dashed,black!40] (0,0) circle (1.35);
\foreach \a/\b in {0.1/0.2, -0.5/0.4, 0.6/-0.3, -0.2/-0.7, 0.9/0.5,
                   -0.9/-0.2, 0.3/0.9, -0.6/-0.8, 1.1/-0.1, -1.0/0.6,
                   0.0/-1.1, 0.5/0.6, -0.3/0.1, 0.75/-0.75, -0.75/0.85,
                   0.2/-0.4, -0.4/-0.3, 1.0/0.2, -1.15/-0.4, 0.45/1.05,
                   -0.15/0.55, 0.65/0.05, -0.55/-0.55, 0.25/-0.9, -0.85/0.25}
  \fill[black!72] (\a,\b) circle (1.5pt);
\draw[dashed,black!55,->] (0,0) -- (1.9,0);
\draw[dashed,black!55,->] (0,0) -- (0,1.9);
\begin{scope}[rotate=32]
\draw[dashed,black!55,->] (0,0) -- (1.9,0);
\draw[dashed,black!55,->] (0,0) -- (0,1.9);
\end{scope}
\node[font=\small] at (0,-2.4) {源高斯：怎么转都一样};
\end{scope}
\end{tikzpicture}
\end{center}

\emph{白化把点云的二阶结构抹平之后，非高斯源留下的方形轮廓仍指明了方向；高斯源的点云旋转对称，方向的信息已经不存在了。}

图左边那个方框说明了 ICA 的全部依据。两个独立的均匀分布混合、白化之后，点云是一个旋转过的正方形，它的边指出了正确的旋转角；沿着任何别的方向去看，边缘分布都会因为多个源叠加而变得更接近钟形——中心极限定理正是往高斯方向推。所以 ICA 的做法是反着走：在所有白化后的方向里，挑边缘分布\emph{最不像高斯}的那些。常用的度量有负熵的近似、峰度，或者等价地最小化各分量之间的互信息 \(I(y_1;\dots;y_k)\)，因为分量独立当且仅当互信息为零。

图右边解释了这条路线的硬边界。若源本身就是高斯的，白化后点云旋转对称，非高斯性对所有方向一样是零，没有任何统计量能区分它们。所以 ICA 的假设必须写成：各分量独立，且\textbf{至多一个是高斯}。眨眼伪迹之所以能被 ICA 抓出来，正因为它极其非高斯——绝大部分时间接近零，偶尔一个大尖峰，峰度很高，与背景脑电的近高斯性质截然不同。

\subsection{ICA 恢复方向，但恢复不了顺序和尺度}

即使假设全部成立，可辨认性也只到某个程度。交换两个源的顺序、同时交换 \(A\) 的对应两列，观测一模一样；把某个源乘以常数 \(c\)、同时把 \(A\) 的对应列除以 \(c\)，观测也一模一样。所以 ICA 至多恢复到置换和尺度：算法输出的``第一个分量''没有``最重要''的含义（与 PCA 不同，那里的顺序由方差定），单个分量的幅度和正负号也是约定出来的，实现里通常靠固定单位方差和某个符号规则钉住。

假设不成立时的失效更值得留意。两个以上的源都是高斯，方向就不可辨认，算法照样会输出一组分量，只是它们与真实源无关。源之间存在时间依赖或本身不独立时，也需要换用利用时间结构或非平稳性的变体。因此，把某个分离出来的波形命名为``眨眼''之前，至少要做两件事：换随机种子和不同的数据段重跑，看该分量是否稳定复现；把它减掉之后检查重构残差和已知的生理特征（眨眼的头皮分布应集中在前额电极）。分离算法给出的是一组坐标，命名是你加上去的。

\subsection{CCA 换一个可辨认的来源：跨视图的相关}

ICA 从单个视图的高阶结构里找方向。另一条路是根本不动高阶统计量，改为引入第二个视图：如果同一批对象上有两组测量，那么``两边一起变''本身就能定住方向。

设成对观测 \(X\in\R^p\) 与 \(Y\in\R^q\)，比如同一天的消费记录与出行记录，或同一批被试的行为指标与影像指标。典型相关分析找投影方向 \(a,b\)，使

\[
\Corr\bigl(a^{\trans}X,\ b^{\trans}Y\bigr)
\]

最大。中心化后，第一对方向是

\[
\max_{a,b}\ a^{\trans}\Sigma_{XY}b
\qquad\text{s.t.}\qquad
a^{\trans}\Sigma_{XX}a=1,\quad b^{\trans}\Sigma_{YY}b=1 ,
\]

两条约束把各自的任意缩放去掉。把两个视图分别白化之后，这个问题化归为对

\[
\Sigma_{XX}^{-1/2}\,\Sigma_{XY}\,\Sigma_{YY}^{-1/2}
\]

做 SVD，奇异值就是典型相关系数，左右奇异向量给出各自的方向。绕了一圈又回到第二节的工具——本单元几乎所有方法最后都落在某个矩阵的谱上，区别只在于对哪个矩阵取谱。

\subsection{高维下的求逆会造出漂亮而不可复现的相关}

CCA 的公式里有两个逆平方根，麻烦就在这儿。当 \(p\) 或 \(q\) 接近样本数时，样本协方差接近奇异，\(\Sigma_{XX}^{-1/2}\) 会把小特征值方向放大到失控。后果是 CCA 几乎必然找到典型相关接近 1 的方向：\(p+q\) 个自由度去拟合 \(n\) 个样本，总有一组系数能让两边对齐。这样的方向在新样本上通常一文不值。

标准补救是正则化，把 \(\Sigma_{XX}+\lambda_x I\) 和 \(\Sigma_{YY}+\lambda_y I\) 代进去做白化，\(\lambda\) 只能在训练折内选，典型相关必须在独立样本上重新计算才算数。这与第一节 PCA 的纪律是同一条，见\hyperref[sec:11-Mathematical-Statistics/06-Resampling-and-Model-Selection/03]{``交叉验证评估的是完整训练流程''}一节。另外，样本协方差在高维下谱的系统性失真，本身就是判断``这个相关有多少是真的''的背景知识，见\hyperref[ch:094]{``随机矩阵理论：高维协方差的谱失真''}一章。

还有一层与统计无关的限制。CCA 找到的是两个视图之间最强的线性共同变化，它不区分这份共同变化从哪来。若日期、季节、采样设备同时影响两边，典型相关会忠实地把这个混杂捕捉下来，并且系数很漂亮。要把它读成共享的因果因素，需要的是设计和识别假设，不是更大的相关系数。

\subsection{同一个单元里，七种方法丢掉的是不同的东西}

回头看整个单元，会发现所有方法都在做同一个动作的不同版本：声明哪些方向的信息可以丢。PCA 丢方差小的方向，代价是可能连判别方向一起丢；截断 SVD 丢强度弱的通道，最优性由 Eckart--Young 担保，而担保在加上非负、稀疏、量化、掩码之后立刻失效；矩阵补全丢的是低秩之外的自由度，能不能恢复取决于观测怎么落；流形方法丢的是全局距离、保住的是局部邻域，于是图上的簇间距离和簇大小都不能读；随机投影几乎不挑方向，只承诺点对距离；因子分析丢的是特有噪声，可辨认的只有载荷子空间；ICA 和 CCA 则各押一条额外假设——非高斯性、跨视图相关——把子空间里的具体方向重新钉住。

这份清单没有优劣排序。ICA 不是比 PCA 更高级的替代品，CCA 也不是。它们回答的是不同的可辨认性问题，各自的答案只在各自的假设下有效。选方法之前先问清楚一件事：这次任务允许丢掉什么？答案定了，方法基本也就定了；答案没定就开始调参，调出来的东西没有办法验证。
